\section{Visitor}
\label{sec:pat-visitor}

\problemstatement{Represent an operation to be performed on the elements of
an object structure, letting you define a new operation without changing
the classes of the elements it operates on.}

\begin{patternbox}[When You Reach for It]
The trigger is \emph{``add new operations over a stable set of element
types''} -- an AST you want to pretty-print, type-check, and compile,
without editing every node class each time. When element types rarely
change but operations keep being added, Visitor centralises each operation.
\end{patternbox}

\subsection*{Understanding the pattern}

Each element exposes \texttt{accept(visitor)}, which calls back
\texttt{visitor.visitConcreteType(*this)} -- a technique called
\textbf{double dispatch}, because the executed method depends on both the
element's type and the visitor's type. A visitor bundles one operation
across all element types as a set of \texttt{visit} overloads. Adding an
operation is a new visitor class; the element classes are untouched.

\begin{ideabox}[Double Dispatch Moves Operations Out of the Elements]
\texttt{accept} resolves the element's concrete type; the matching
\texttt{visit} overload resolves the operation. Together they route to the
right (element, operation) pair without a type switch, so operations live
in visitors rather than being smeared across every element class.
\end{ideabox}

\subsection*{C++ implementation}

\begin{lstlisting}
#include <bits/stdc++.h>
using namespace std;

struct Circle; struct Square;
struct Visitor {
    virtual void visit(Circle&) = 0;
    virtual void visit(Square&) = 0;
    virtual ~Visitor() = default;
};

struct Shape { virtual void accept(Visitor& v) = 0; virtual ~Shape() = default; };
struct Circle : Shape {
    double r = 2;
    void accept(Visitor& v) override { v.visit(*this); }   // double dispatch
};
struct Square : Shape {
    double side = 3;
    void accept(Visitor& v) override { v.visit(*this); }
};

struct AreaVisitor : Visitor {                  // a new operation, no shape edits
    void visit(Circle& c) override { cout << "Circle area " << 3.14159 * c.r * c.r << "\n"; }
    void visit(Square& s) override { cout << "Square area " << s.side * s.side << "\n"; }
};

int main() {
    vector<unique_ptr<Shape>> shapes;
    shapes.push_back(make_unique<Circle>());
    shapes.push_back(make_unique<Square>());
    AreaVisitor area;
    for (auto& s : shapes) s->accept(area);     // right visit() per type
    return 0;
}
\end{lstlisting}

\begin{complexitybox}[Trade-offs]
\textbf{Pros.} Add new operations easily (one new visitor); gathers a
scattered operation into one class; keeps element classes clean.
\textbf{Cons.} Adding a new \emph{element type} forces editing every
visitor -- the exact opposite trade-off to most patterns. Use only when
element types are stable and operations grow.
\end{complexitybox}

\begin{takeawaybox}
Visitor externalises operations over an object structure via double
dispatch, so new operations are new visitor classes with the elements
untouched -- at the cost of making new element types expensive. It is the
inverse trade-off of a type switch, and the go-to for AST/compiler passes.
\end{takeawaybox}
